\begin{table}[htbp]
\centering
\caption{Standardized Effect Sizes}
\label{tab:sde}
\footnotesize
\begin{tabular}{lcccccc}
\hline\hline
Outcome & $\hat{\beta}$ & SE & SD($Y$) & SDE & SE(SDE) & Classification \\
\hline
\multicolumn{7}{l}{\textit{Panel A: Pooled}} \\
Violation rate & 0.0030 & 0.0163 & 0.3195 & 0.0095 & 0.0509 & Small positive \\
Any violation & 0.0021 & 0.0187 & 0.3472 & 0.0062 & 0.0540 & Small positive \\
\hline
\multicolumn{7}{l}{\textit{Panel B: Heterogeneous (by watershed violation intensity)}} \\
High-violation watersheds & -0.0318 & 0.0200 & 0.2900 & -0.1097 & 0.0691 & Moderate negative \\
Low-violation watersheds & 0.0336 & 0.0227 & 0.2456 & 0.1367 & 0.0925 & Moderate positive \\
\hline\hline
\end{tabular}
\begin{tablenotes}[flushleft]\footnotesize
\item \textit{Notes:} \textbf{Country:} United States. \textbf{Research question:} Does Clean Water Act Section 303(d) impaired-waters listing cause permitted facilities to reduce compliance violations? \textbf{Policy mechanism:} 303(d) listing designates a water body as impaired under the CWA, triggering Total Maximum Daily Load (TMDL) development and wasteload allocations that impose stricter effluent limits on point-source dischargers; facilities in listed watersheds face tighter permits and increased EPA/state enforcement scrutiny. \textbf{Outcome definition:} Share of 13 most recent quarters with any compliance violation (effluent exceedance, significant noncompliance, or permit violation) from EPA ECHO quarterly compliance tracking. \textbf{Treatment:} Binary --- facility's HUC-12 subwatershed is on the state's 303(d) list of impaired waters. \textbf{Data:} EPA ECHO facility compliance data and ATTAINS 303(d) impairment assessments, 2024 reporting cycle, facility-level cross-section, 2,276 NPDES-permitted facilities. \textbf{Method:} OLS with HUC-8 watershed fixed effects; standard errors clustered at HUC-8 level. \textbf{Sample:} Boundary sample restricted to HUC-8 watersheds containing both 303(d)-listed and non-listed HUC-12 subwatersheds, ensuring within-watershed comparisons. SDE $= \hat{\beta} / \text{SD}(Y)$ where SD($Y$) is the control-group standard deviation. Classification refers to magnitude, not statistical significance: Large ($|$SDE$| > 0.15$), Moderate ($0.05$--$0.15$), Small ($0.005$--$0.05$), Null ($< 0.005$). 
\end{tablenotes}
\end{table}
